% Seccion 5 de la rubrica: modelacion y validacion. Solo abre la seccion y
% fija lo comun a todas las tecnicas (particion, regla de fuga). Cada tecnica
% es una \subsection en su propio archivo.

\section{Modelación y validación}
\label{sec:modelacion}

Todas las técnicas comparten la misma partición temporal: los años
2021--2024 se destinan al ajuste y 2025 completo a la validación. La regla
es estricta: ningún procedimiento ve 2025 durante el ajuste, incluidos el
ACP y el AF, cuyos ejes se estiman sobre 2021--2024 y se aplican a 2025 con
los pesos de entrenamiento; de otro modo la información de la partición de
validación se filtraría a las coordenadas que consumen los clasificadores.
La única fuga que persiste es la del parámetro $\lambda$ de Box-Cox,
estimado sobre el periodo completo en \cref{sec:transformaciones}; se
declara porque revertirla exigiría rehacer el conjunto transformado.

La validación tiene una particularidad que atraviesa todos los resultados:
2025 es el año de ozono más alto de la serie, con 44.0\,\% de excedencia
día-estación contra 29.4\,\% del ajuste, en las cuatro temporadas. Es un
efecto de año, no de calendario, y cada técnica lo retoma donde le afecta.

Las cinco técnicas se presentan en el orden en que cada una alimenta a la
siguiente: el ACP (\cref{sec:pca}) y el AF (\cref{sec:analisis-factorial})
reducen los 16 predictores a tres ejes; la caracterización como serie de
tiempo (\cref{sec:serie-tiempo}) establece la dependencia temporal de esos
ejes; la regresión logística y el discriminante
(\cref{sec:clasificacion_tecs}) clasifican el día-estación; y el modelo
autorregresivo (\cref{sec:autorregresivo}) incorpora la dependencia temporal
a la clasificación.
